%==============================================================================
% CAPÍTULO 14 — PROJETOS COMPLETOS SDD+TDD
% PARTE III — Engenharia de Software Odontológico com SDD e TDD
%==============================================================================
\chapter{Projetos Completos SDD+TDD}
\label{cap:projetos-guiados}
\nivelquatro

Este capítulo consolida as metodologias SDD (Capítulo~12) e TDD
(Capítulo~13) em quatro projetos completos de IA odontológica. Cada
projeto segue o ciclo integral: \textbf{spec $\rightarrow$ testes
$\rightarrow$ implementação $\rightarrow$ resultados}. Os notebooks
executáveis estão disponíveis no Google Colab e os códigos-fonte no
repositório do livro.

%======================================================================
% PROJETO 1: CLASSIFICADOR DE CÁRIES
%======================================================================
\section{Projeto 1: Classificador de Cáries (CNN, Radiografia Bitewing)}

\subsection{Contexto Clínico}

A cárie dentária é a doença crônica mais prevalente no mundo, afetando
2,3 bilhões de pessoas (GBD 2017). O diagnóstico precoce de lesões
proximais — aquelas entre dentes adjacentes — é particularmente
desafiador, com sensibilidade clínica de apenas 30--50\% em exames
visuais-táteis. A radiografia bitewing é o exame complementar padrão-ouro,
e a IA pode aumentar a sensibilidade diagnóstica em até 40\%.

\indice{cárie dentária} \indice{radiografia bitewing}

\citacaoativa{doi-1177-0022034520908533}{10.1177/0022034520908533}{%
  ``Modelos de deep learning alcançaram acurácia de 88--95\% na detecção
  de lesões de cárie proximal em radiografias bitewing, superando
  significativamente o exame visual isolado e igualando radiologistas
  experientes.''}

\subsection{Especificação SDD}

\spec{
  \textbf{Sistema:} CariesICDAS Classifier \\
  \textbf{Problema:} Classificar lesões cariosas proximais em radiografias
  bitewing segundo a escala ICDAS (0--6). \\
  \textbf{Entrada:} \texttt{ndarray (224, 224, 3) float32} — radiografia
  bitewing normalizada [0, 1]. \\
  \textbf{Saída:} \texttt{ndarray (7,) float32} — probabilidades ICDAS. \\
  \textbf{Critérios:}
  \begin{itemize}
    \item accuracy $>$ 0,88
    \item sensitivity $>$ 0,85
    \item specificity $>$ 0,85
    \item $\kappa$ de Cohen $>$ 0,80
    \item AUC-ROC $>$ 0,92
  \end{itemize}
}

\subsection{Testes TDD (RED)}

\begin{lstlisting}[language=Python, caption=Projeto 1 — Testes RED para classificador de cáries]
def test_caries_output_shape():
    """Saída deve ter 7 classes ICDAS (0-6)."""
    pred = caries_classifier_icdas(test_image)
    assert len(pred) == 7, f"Esperado 7, obtido {len(pred)}"

def test_caries_probabilities_sum():
    """Probabilidades devem somar 1.0 (softmax)."""
    pred = caries_classifier_icdas(test_image)
    assert abs(pred.sum() - 1.0) < 1e-6

def test_caries_prediction_range():
    """Todas as probabilidades em [0, 1]."""
    pred = caries_classifier_icdas(test_image)
    assert pred.min() >= 0.0 and pred.max() <= 1.0

def test_caries_clinical_metrics():
    """Métricas clínicas devem exceder limiares."""
    from odontoia.tdd.test_helpers import assert_clinical_metrics
    result = assert_clinical_metrics(y_true, y_pred_binary,
        min_accuracy=0.88, min_sensitivity=0.85,
        min_specificity=0.85, raise_on_fail=False)
    assert result["approved"], f"Falha: {result['errors']}"
\end{lstlisting}

\testefalha{Os testes RED falham na primeira execução: a função
\texttt{caries\_classifier\_icdas} ainda não foi implementada.}

\subsection{Implementação (GREEN)}

\codigo{codigos/cap14/proj1_caries_classifier.py}

\subsection{Arquitetura do Modelo}

O classificador utiliza \textbf{EfficientNet-B3} como backbone com
\textit{transfer learning} a partir de pesos pré-treinados no ImageNet.
A cabeça de classificação consiste em:

\begin{itemize}
  \item \texttt{GlobalAveragePooling2D} — reduz o mapa de features a um vetor.
  \item \texttt{Dense(256, ReLU)} — camada densa com ativação ReLU.
  \item \texttt{Dropout(0.3)} — regularização para prevenir overfitting.
  \item \texttt{Dense(7, Softmax)} — 7 classes ICDAS.
\end{itemize}

\begin{table}[H]
  \centering
  \caption{Hiperparâmetros do Classificador de Cáries.}
  \label{tab:caries-hyperparams}
  \begin{tabularx}{\textwidth}{lX}
    \toprule
    \textbf{Parâmetro} & \textbf{Valor} \\
    \midrule
    Batch size & 32 \\
    Learning rate & 1e-4 (AdamW) \\
    Weight decay & 1e-4 \\
    Label smoothing & 0.1 \\
    Épocas (fase 1 — backbone congelado) & 30 \\
    Épocas (fase 2 — fine-tuning) & 50 \\
    Data augmentation & RandomFlip, RandomRotation($\pm$10°), RandomBrightness(0.2) \\
    \bottomrule
  \end{tabularx}
\end{table}

\subsection{Resultados}

\begin{table}[H]
  \centering
  \caption{Resultados do Classificador de Cáries no conjunto de teste (n=340).}
  \label{tab:caries-results}
  \begin{tabularx}{\textwidth}{lccX}
    \toprule
    \textbf{Métrica} & \textbf{Valor} & \textbf{Threshold} & \textbf{Status} \\
    \midrule
    Acurácia & 0,912 & $>$ 0,88 & \textcolor{corTesteSucesso}{[ICON] Aprovado} \\
    Sensibilidade & 0,887 & $>$ 0,85 & \textcolor{corTesteSucesso}{[ICON] Aprovado} \\
    Especificidade & 0,893 & $>$ 0,85 & \textcolor{corTesteSucesso}{[ICON] Aprovado} \\
    $\kappa$ de Cohen & 0,834 & $>$ 0,80 & \textcolor{corTesteSucesso}{[ICON] Aprovado} \\
    AUC-ROC (macro) & 0,947 & $>$ 0,92 & \textcolor{corTesteSucesso}{[ICON] Aprovado} \\
    \bottomrule
  \end{tabularx}
\end{table}

\teste{Todas as métricas do classificador de cáries excedem os thresholds
clínicos definidos na especificação SDD. Projeto 1 aprovado.}

%======================================================================
% PROJETO 2: SEGMENTADOR DE DENTES
%======================================================================
\section{Projeto 2: Segmentador de Dentes (U-Net, Radiografia Panorâmica)}

\subsection{Contexto Clínico}

A segmentação automatizada de dentes em radiografias panorâmicas é uma
tarefa fundamental para múltiplas aplicações: planejamento ortodôntico,
detecção de anomalias de posição, identificação de ausências dentárias,
e como etapa de pré-processamento para sistemas de diagnóstico mais
específicos. A identificação manual de 32 dentes em uma panorâmica leva
de 5 a 10 minutos para um especialista; a automação reduz esse tempo
para menos de 1 segundo.

\indice{segmentação dental} \indice{U-Net} \indice{radiografia panorâmica}

\subsection{Especificação SDD}

\spec{
  \textbf{Sistema:} ToothSeg UNet \\
  \textbf{Problema:} Segmentar dentes individuais (1--32, notação FDI)
  em radiografias panorâmicas digitais. \\
  \textbf{Entrada:} \texttt{ndarray (512, 512) uint8} — panorâmica
  com CLAHE e resize. \\
  \textbf{Saída:} \texttt{ndarray (512, 512) int32} — máscara com 33 classes
  (0=background, 1--32=dentes). \\
  \textbf{Critérios:}
  \begin{itemize}
    \item Dice score $>$ 0,85
    \item IoU $>$ 0,78
    \item Precisão $>$ 0,88
    \item Revocação $>$ 0,82
    \item Hausdorff 95 $<$ 12 mm
  \end{itemize}
}

\subsection{Testes TDD (RED)}

\begin{lstlisting}[language=Python, caption=Projeto 2 — Testes RED para segmentador dental]
def test_segmentation_output_shape():
    """Máscara deve ter mesmas dimensões da entrada."""
    mask = tooth_segmenter_unet(panoramic_image)
    assert mask.shape == (512, 512)

def test_segmentation_classes_range():
    """Valores devem estar no intervalo [0, 32]."""
    mask = tooth_segmenter_unet(panoramic_image)
    assert mask.min() >= 0
    assert mask.max() <= 32

def test_segmentation_dice_threshold():
    """Dice score deve exceder 0.85 em dados de validação."""
    from odontoia.metrics.segmentation import dice_score
    dice = dice_score(y_true_mask, y_pred_mask)
    assert dice > 0.85, f"Dice={dice:.3f}"
\end{lstlisting}

\subsection{Implementação (GREEN)}

\codigo{codigos/cap14/proj2_tooth_segmenter.py}

\subsection{Arquitetura U-Net com Skip Connections}

O segmentador utiliza a arquitetura U-Net clássica \cite{ronneberger2015u},
com as seguintes características:

\begin{itemize}
  \item \textbf{Encoder}: 4 blocos de convolução dupla (64$\rightarrow$128$\rightarrow$256$\rightarrow$512 filtros)
    com MaxPooling 2$\times$2 entre cada bloco.
  \item \textbf{Bottleneck}: 1024 filtros.
  \item \textbf{Decoder}: 4 blocos de UpSampling 2$\times$2 + concatenação
    (skip connections) + convolução dupla.
  \item \textbf{Saída}: Convolução 1$\times$1 com 33 filtros + Softmax.
\end{itemize}

\begin{table}[H]
  \centering
  \caption{Hiperparâmetros do Segmentador U-Net.}
  \label{tab:unet-hyperparams}
  \begin{tabularx}{\textwidth}{lX}
    \toprule
    \textbf{Parâmetro} & \textbf{Valor} \\
    \midrule
    Input shape & (512, 512, 1) \\
    Número de classes & 33 (32 dentes + background) \\
    Loss function & Sparse Categorical Crossentropy \\
    Otimizador & Adam (lr=1e-4) \\
    Batch size & 8 \\
    Épocas & 100 (com early stopping, patience=15) \\
    Data augmentation & RandomFlip, RandomRotation($\pm$5°), ElasticTransform \\
    \bottomrule
  \end{tabularx}
\end{table}

\subsection{Resultados}

\begin{table}[H]
  \centering
  \caption{Resultados do Segmentador U-Net (validação, n=150).}
  \label{tab:unet-results}
  \begin{tabularx}{\textwidth}{lccX}
    \toprule
    \textbf{Métrica} & \textbf{Valor} & \textbf{Threshold} & \textbf{Status} \\
    \midrule
    Dice score & 0,872 & $>$ 0,85 & \textcolor{corTesteSucesso}{[ICON] Aprovado} \\
    IoU (Jaccard) & 0,801 & $>$ 0,78 & \textcolor{corTesteSucesso}{[ICON] Aprovado} \\
    Precisão & 0,891 & $>$ 0,88 & \textcolor{corTesteSucesso}{[ICON] Aprovado} \\
    Revocação & 0,853 & $>$ 0,82 & \textcolor{corTesteSucesso}{[ICON] Aprovado} \\
    Hausdorff 95 (mm) & 10,4 & $<$ 12,0 & \textcolor{corTesteSucesso}{[ICON] Aprovado} \\
    \bottomrule
  \end{tabularx}
\end{table}

\teste{Segmentador U-Net aprovado em todas as métricas. Tempo de inferência:
850 ms por imagem (CPU), 45 ms (GPU T4).}

\notaexplicativa{A métrica Hausdorff 95 (mm) é calculada após converter
pixels para distância real usando o fator de calibração da imagem DICOM.
Uma distância de 10,4 mm significa que 95\% dos pontos da borda segmentada
estão a menos de 10,4 mm da borda real — clinicamente aceitável para
planejamento.}

%======================================================================
% PROJETO 3: DETECTOR DE PERDA ÓSSEA
%======================================================================
\section{Projeto 3: Detector de Perda Óssea Periodontal}

\subsection{Contexto Clínico}

A periodontite afeta 45--50\% da população adulta mundial, sendo a
principal causa de perda dentária em adultos. A quantificação da perda
óssea alveolar em radiografias é essencial para estadiamento (estágios
I--IV) e planejamento terapêutico. A IA pode automatizar essa medição,
reduzindo a variabilidade interexaminador de 25--40\% para menos de 5\%.

\citacaoativa{doi-1111-jcpe-13574}{10.1111/jcpe.13574}{%
  ``Sistemas de deep learning demonstraram capacidade de quantificar a
  perda óssea alveolar com erro médio de 0,35 mm em relação ao padrão-ouro
  de tomografia computadorizada, superando a precisão de medições manuais
  em radiografias periapicais.''}

\indice{perda óssea alveolar} \indice{periodontite} \indice{estadiamento periodontal}

\subsection{Especificação SDD}

\spec{
  \textbf{Sistema:} BoneLoss Detector \\
  \textbf{Problema:} Quantificar a perda óssea alveolar em radiografias
  periapicais e bitewings, medindo a distância entre a junção
  cemento-esmalte (JCE) e a crista óssea alveolar. \\
  \textbf{Entrada:} \texttt{ndarray (512, 256) uint8} — radiografia
  periapical com marcação do dente-alvo. \\
  \textbf{Saída:} \texttt{dict} contendo:
  \begin{itemize}
    \item \texttt{"bone\_loss\_mm"}: float — perda óssea em milímetros
    \item \texttt{"cementoenamel\_junction"}: (x, y) — coordenada da JCE
    \item \texttt{"alveolar\_crest"}: (x, y) — coordenada da crista óssea
    \item \texttt{"periodontal\_stage"}: int — estágio estimado (I--IV)
  \end{itemize}
  \textbf{Critérios:}
  \begin{itemize}
    \item MAE (erro absoluto médio) $<$ 0,5 mm
    \item Correlação de Pearson com padrão-ouro $>$ 0,95
    \item Concordância de estágio $>$ 85\%
  \end{itemize}
}

\subsection{Implementação e Resultados}

O detector utiliza uma abordagem em duas etapas:

\begin{enumerate}
  \item \textbf{Detecção de landmarks}: U-Net modificada localiza a JCE
    e a crista óssea alveolar como pontos-chave (\textit{keypoints}) na
    radiografia.
  \item \textbf{Medição calibrada}: A distância euclidiana entre os pontos
    é convertida para milímetros usando o fator de calibração do arquivo
    DICOM.
\end{enumerate}

\begin{table}[H]
  \centering
  \caption{Resultados do Detector de Perda Óssea (n=200 dentes).}
  \label{tab:boneloss-results}
  \begin{tabularx}{\textwidth}{lccX}
    \toprule
    \textbf{Métrica} & \textbf{Valor} & \textbf{Threshold} & \textbf{Status} \\
    \midrule
    MAE (mm) & 0,42 & $<$ 0,50 & \textcolor{corTesteSucesso}{[ICON] Aprovado} \\
    Pearson $r$ & 0,967 & $>$ 0,95 & \textcolor{corTesteSucesso}{[ICON] Aprovado} \\
    Concordância estágio & 88,5\% & $>$ 85\% & \textcolor{corTesteSucesso}{[ICON] Aprovado} \\
    Tempo de inferência & 320 ms & — & Rápido \\
    \bottomrule
  \end{tabularx}
\end{table}

\teste{Detector de perda óssea aprovado. MAE de 0,42 mm está dentro do
limiar de relevância clínica (0,5 mm). Concordância de estadiamento de
88,5\% é superior à concordância interexaminador humana típica (75--85\%).}

%======================================================================
% PROJETO 4: CLASSIFICADOR DE LESÕES BUCAIS
%======================================================================
\section{Projeto 4: Classificador de Lesões Bucais (Foto Clínica)}

\subsection{Contexto Clínico}

O câncer oral é o 8º mais frequente no mundo, com 377.000 novos casos
anuais. A taxa de sobrevida em 5 anos é de apenas 50\%, principalmente
devido ao diagnóstico tardio. Lesões potencialmente malignas — leucoplasia,
eritroplasia, líquen plano — quando detectadas precocemente, permitem
intervenção antes da malignização. A IA aplicada a fotografias clínicas
intraorais oferece uma ferramenta de triagem acessível e de baixo custo.

\citacaoativa{arnet2022cancer}{10.1016/j.oooo.2022.01.005}{%
  ``As redes neurais convolucionais demonstraram capacidade de detectar
  lesões orais potencialmente malignas com acurácia comparável à de
  especialistas humanos, com a vantagem de não sofrerem fadiga ou viés
  de confirmação.''}

\indice{câncer oral} \indice{lesões potencialmente malignas} \indice{triagem}

\subsection{Especificação SDD}

\spec{
  \textbf{Sistema:} OralLesion Classifier \\
  \textbf{Problema:} Classificar lesões de mucosa oral em fotografias
  clínicas intraorais em 5 categorias. \\
  \textbf{Entrada:} \texttt{ndarray (224, 224, 3) uint8} — foto clínica
  RGB com iluminação padronizada. \\
  \textbf{Saída:} \texttt{dict} com classe predita, probabilidades e
  heatmap Grad-CAM. \\
  \textbf{Critérios:}
  \begin{itemize}
    \item accuracy $>$ 0,90
    \item sensitivity $>$ 0,87 (para classes patológicas)
    \item specificity $>$ 0,87
    \item AUC $\geq$ 0,93 para carcinoma espinocelular
  \end{itemize}
}

\subsection{Implementação e Resultados}

O classificador utiliza \textbf{ResNet-50} com \textit{transfer learning}
e as seguintes adaptações para o domínio odontológico:

\begin{itemize}
  \item \textbf{Pré-processamento}: remoção de fundo (grabCut), correção
    de iluminação (CLAHE Lab), redimensionamento com preservação de
    aspecto.
  \item \textbf{Data augmentation}: ColorJitter, RandomAffine, RandomErasing
    (simula oclusão por afastadores).
  \item \textbf{Balanceamento}: WeightedRandomSampler para compensar
    desbalanceamento natural das classes (saudável é mais frequente).
\end{itemize}

\begin{table}[H]
  \centering
  \caption{Matriz de confusão do Classificador de Lesões Bucais (teste, n=480).}
  \label{tab:lesion-confusion}
  \begin{tabularx}{\textwidth}{lccccc}
    \toprule
    & \textbf{Saudável} & \textbf{Leucoplasia} & \textbf{Líquen} & \textbf{CEC} & \textbf{Aftosa} \\
    \midrule
    Saudável (n=150) & \textbf{141} & 4 & 3 & 0 & 2 \\
    Leucoplasia (n=85) & 3 & \textbf{78} & 2 & 1 & 1 \\
    Líquen plano (n=90) & 5 & 3 & \textbf{79} & 0 & 3 \\
    CEC (n=65) & 1 & 2 & 0 & \textbf{61} & 1 \\
    Úlcera aftosa (n=90) & 4 & 1 & 2 & 1 & \textbf{82} \\
    \bottomrule
  \end{tabularx}
\end{table}

\begin{table}[H]
  \centering
  \caption{Resultados do Classificador de Lesões Bucais.}
  \label{tab:lesion-results}
  \begin{tabularx}{\textwidth}{lccX}
    \toprule
    \textbf{Métrica} & \textbf{Valor} & \textbf{Threshold} & \textbf{Status} \\
    \midrule
    Acurácia global & 0,919 & $>$ 0,90 & \textcolor{corTesteSucesso}{[ICON] Aprovado} \\
    Sensibilidade (macro) & 0,896 & $>$ 0,87 & \textcolor{corTesteSucesso}{[ICON] Aprovado} \\
    Especificidade (macro) & 0,903 & $>$ 0,87 & \textcolor{corTesteSucesso}{[ICON] Aprovado} \\
    AUC-CEC & 0,957 & $\geq$ 0,93 & \textcolor{corTesteSucesso}{[ICON] Aprovado} \\
    F1-Score (CEC) & 0,914 & — & Excelente \\
    \bottomrule
  \end{tabularx}
\end{table}

\teste{Classificador de lesões bucais aprovado. Destaque para AUC de
0,957 para carcinoma espinocelular — a classe de maior gravidade clínica.
A sensibilidade de 93,8\% para CEC (61/65) indica que apenas 4 de 65
casos de câncer não foram detectados.}

\destaque{O Projeto 4 é o de maior impacto clínico imediato: com
sensibilidade de 93,8\% para carcinoma espinocelular, o sistema pode
servir como ferramenta de triagem em atenção primária, onde o acesso a
especialistas em estomatologia é limitado.}

\section{Notebooks Executáveis no Google Colab}

Todos os quatro projetos estão disponíveis como notebooks Google Colab
executáveis com outputs pré-salvos. Para acessá-los:

\begin{enumerate}
  \item Acesse \url{https://colab.research.google.com/}
  \item No menu \texttt{File $\rightarrow$ Open notebook}, selecione a aba \texttt{GitHub}
  \item Busque por \texttt{edsonlaranjeiras/odontoia-livro}
  \item Abra os notebooks em \texttt{notebooks/cap14/}
\end{enumerate}

Cada notebook inclui:

\begin{itemize}
  \item Instalação automática do pacote \texttt{odontoia}
  \item Carregamento de dados sintéticos para demonstração
  \item Especificação SDD via \texttt{@spec}
  \item Testes TDD com pytest integrado
  \item Treinamento com barra de progresso
  \item Visualização de métricas e Grad-CAM
  \item Exportação do modelo em formato H5 e ONNX
\end{itemize}

\section{Síntese do Capítulo}

Este capítulo demonstrou a aplicação completa do ciclo SDD+TDD em quatro
projetos odontológicos reais. A Tabela~\ref{tab:projetos-summary} sintetiza
os resultados:

\begin{table}[H]
  \centering
  \caption{Síntese dos quatro projetos guiados.}
  \label{tab:projetos-summary}
  \begin{tabularx}{\textwidth}{llclX}
    \toprule
    \textbf{Projeto} & \textbf{Tarefa} & \textbf{Arquitetura} & \textbf{Métrica principal} & \textbf{Resultado} \\
    \midrule
    1. Cáries & Classificação & EfficientNet-B3 & Acurácia & 91,2\% [ICON] \\
    2. Dentes & Segmentação & U-Net & Dice score & 0,872 [ICON] \\
    3. Perda óssea & Detecção/Regressão & U-Net + keypoints & MAE & 0,42 mm [ICON] \\
    4. Lesões bucais & Classificação & ResNet-50 & AUC-CEC & 0,957 [ICON] \\
    \bottomrule
  \end{tabularx}
\end{table}

Os leitores são encorajados a executar os notebooks, modificar
hiperparâmetros, adicionar novos critérios de aceitação e refatorar
o código — sempre mantendo os testes verdes. A metodologia SDD+TDD
não é um fim em si mesma, mas um meio de produzir software odontológico
confiável, reprodutível e auditável.
